\begin{helper}
Fix a finite index set \(\Omega\), a finite chain length \(n\), a cell-weight
function \(w:\Omega\to\mathbb R\), endpoint representatives
\(S_\omega,T_\omega\) indexed by \(\omega\in\Omega\), and adjacent
representatives \(L_{k,\omega},R_{k,\omega}\) indexed by \(k<n\) and
\(\omega\in\Omega\).  Assume all these representatives are nonnegative on the
indices where they are used.  Assume also the zero-weight convention
\[
  w(\omega)=0\Rightarrow
  S_\omega=T_\omega=L_{k,\omega}=R_{k,\omega}=0
\]
for every adjacent index \(k<n\) to which the displayed adjacent
representatives apply.  Assume the one-step
inequalities
\[
  L_{k,\omega}\le R_{k,\omega}\qquad(k<n),
\]
the initial handoff \(S_\omega\le R_{0,\omega}\) when \(0<n\), the interior
handoffs
\[
  R_{k,\omega}\le L_{k+1,\omega}\qquad(k+1<n),
\]
the terminal handoff \(R_{n-1,\omega}\le T_\omega\) when \(0<n\), and in the
degenerate case \(n=0\) the endpoint identity \(S_\omega=T_\omega\).
Then there is one global family \(P(j,\omega)\) with
\[
P(j,\omega)\ge0,\qquad w(\omega)=0\Rightarrow P(j,\omega)=0,
\]
\[
P(0,\omega)=S_\omega,\qquad P(n,\omega)=T_\omega,
\]
and
\[
P(k,\omega)\le P(k+1,\omega)\qquad(k<n).
\]
\end{helper}
\begin{proof}
Define \(P(0,\omega)=S_\omega\).  For \(0<j<n\), define
\(P(j,\omega)=R_{j-1,\omega}\).  For \(j\ge n\), define
\(P(j,\omega)=T_\omega\).  This is a single global family, chosen before any
adjacent comparison is checked.

Nonnegativity follows immediately from the corresponding nonnegativity of the
source endpoint, target endpoint, and target adjacent representatives.  If
\(w(\omega)=0\), the same case distinction and the zero-weight convention for
the representative selected in that case give \(P(j,\omega)=0\).  The endpoint
identities are immediate: \(P(0,\omega)=S_\omega\) by definition, and if
\(n=0\) the assumed endpoint identity gives \(P(n,\omega)=T_\omega\), while if
\(0<n\) the clause \(j\ge n\) gives \(P(n,\omega)=T_\omega\).

It remains to check monotonicity.  If \(n=0\), there is no \(k<n\).  If
\(k=0<n\) and \(1<n\), the first comparison is
\[
S_\omega=P(0,\omega)\le R_{0,\omega}=P(1,\omega),
\]
using the initial handoff.  If \(k=0\) and \(n=1\), then
\[
P(0,\omega)=S_\omega\le R_{0,\omega}\le T_\omega=P(1,\omega),
\]
using the initial and terminal handoffs.  For \(0<k\) and \(k+1<n\), the
definitions give
\[
P(k,\omega)=R_{k-1,\omega}\le L_{k,\omega}\le R_{k,\omega}=P(k+1,\omega),
\]
where the first inequality is the interior handoff and the second is the
one-step inequality.  Finally, if \(0<k<n\) and \(k+1=n\), then
\[
P(k,\omega)=R_{k-1,\omega}\le L_{k,\omega}\le R_{k,\omega}\le T_\omega
  =P(n,\omega),
\]
using the interior handoff when \(k>0\), the one-step inequality at \(k\),
and the terminal handoff.  These cases exhaust all \(k<n\), so \(P\) is a
monotone global chain.
\end{proof}
